%% rs_token_v0.4.tex -- RS-Token manuscript with v0.4 experiments
%% Compile: pdflatex rs_token_v0.4.tex; bibtex rs_token_v0.4; pdflatex rs_token_v0.4.tex; pdflatex rs_token_v0.4.tex

\documentclass[lettersize,journal]{IEEEtran}

\usepackage[utf8]{inputenc}
\usepackage{amsmath,amssymb,amsfonts}
\usepackage{graphicx}
\usepackage{textcomp}
\usepackage{xcolor}
\usepackage{booktabs}
\usepackage{multirow}
\usepackage{hyperref}
\usepackage{cite}
\usepackage{url}

\hyphenation{op-tical net-works semi-conduc-tor IEEE-Xplore}

\begin{document}

\title{RS-Token: Hierarchical RemoteCLIP-Distilled Tokens for Channel-Robust Remote Sensing Communication}

\author{Baohui~Zhang%
\thanks{The author is with the School of XXX, Anhui University, Hefei 230601, China (e-mail: y125211005@ahu.edu.cn).}%
}

\markboth{IEEE Geoscience and Remote Sensing Letters,~Vol.~XX, No.~XX, June~2026}%
{Zhang \MakeLowercase{\textit{et al.}}: RS-Token for Remote Sensing Communication}

\maketitle

\begin{abstract}
Remote-sensing image downlinks often operate under limited bandwidth, low SNR, and fading, while receivers may require different levels of information: a scene label for urgent screening, a task-usable low-rate image, or a more faithful reconstruction for visual inspection. This paper presents \textbf{RS-Token}, a hierarchical discrete-token representation for remote-sensing communication. RS-Token uses a four-stage residual vector quantizer (RVQ) and distills RemoteCLIP features into only the first quantization layer, encouraging L0 to carry land-cover semantics while later layers refine image details. The receiver can transmit L0 at 2{,}560 bits/image for the low-rate task path, or add L1--L3 at 5{,}120, 7{,}680, and 10{,}240 bits/image for progressive reconstruction under the present AID configuration. Our evaluation strictly separates the task path from the reconstruction path: L0 bag-of-words accuracy is used for $k=1$ task-path claims, whereas PSNR, LPIPS, and a clean AID ResNet34 reconstructed-image classifier are used for reconstruction claims. Across three model seeds on AID, RemoteCLIP distillation improves no-channel L0 task accuracy from $58.23\pm1.57\%$ to $83.33\pm0.81\%$, and preserves $82.57\pm0.31\%$ under AWGN $+5$ dB and $78.80\pm0.72\%$ under Rayleigh $+10$ dB. Reconstruction-path results show that, in the no-channel sweep, increasing $k$ from 1 to 4 improves PSNR / reconstructed-image classifier accuracy from 22.95 dB / 70.3\% to 25.92 dB / 86.9\%. A layered probe further indicates that L1--L3 add little additional task semantics, supporting the interpretation that they mainly provide residual details. Supplementary unprotected WebP/JPEG2000 and rate-1/2 systematic sparse LDPC experiments clarify external stress-test and channel-coding compatibility boundaries; the LDPC experiment uses a custom sparse code with min-sum belief-propagation decoding and is not a 5G NR LDPC comparison.
\end{abstract}

\begin{IEEEkeywords}
Semantic communication, remote sensing, residual vector quantization, RemoteCLIP, task-oriented communication, hierarchical tokens, channel robustness.
\end{IEEEkeywords}

\section{Introduction}
\IEEEPARstart{R}{emote-sensing} image communication is increasingly used in low-altitude UAV inspection and emergency scene screening. In such settings, the downlink can be unstable because of range variation, occlusion, multipath fading, and temporary communication constraints. At the same time, the receiver does not always need a full-fidelity image immediately. Some real-time alerting tasks only need a reliable scene-level decision; field assessment may need an approximately interpretable image; later mapping, archiving, or manual review may require a higher-quality reconstruction. Remote-sensing communication therefore benefits from a hierarchical representation: at low rate, it should preserve task-relevant semantics, and when more bits can be transmitted, it should progressively recover image details.

Traditional image communication usually follows a compress-then-transmit pipeline. The sender encodes the image into a compressed bitstream, and the receiver decodes the image after recovering the bitstream. This paradigm is mature, but it does not directly match the hierarchical information demand above. On one hand, compressed bitstreams such as WebP or JPEG2000 can be sensitive to uncorrected bit errors, which may cause severe distortion or decoding failure. On the other hand, conventional compression primarily optimizes visual reconstruction and does not explicitly guarantee that the earliest transmitted bits preserve remote-sensing scene semantics. Channel coding can reduce bit errors, but it does not by itself change the source representation's semantic ordering.

A natural alternative is to transmit the discrete indices of a neural tokenizer, such as VQ or RVQ indices. Compared with conventional compressed files, discrete indices are easier to combine with digital channels, retransmission, and error protection, and RVQ naturally supports layer-wise transmission. However, a standard RVQ tokenizer is usually driven by reconstruction losses. Its first codebook mainly approximates a continuous latent at a coarse level and is not guaranteed to be discriminative for remote-sensing scene categories. Thus, plain RVQ index transmission provides a layered rate structure, but it does not necessarily provide low-rate task-semantic fidelity.

RS-Token introduces RemoteCLIP distillation into this layered RVQ representation to explicitly change the role of the first token layer. The distillation loss is applied only to L0-derived features, aligning them with a remote-sensing vision-language embedding. L1--L3 continue to encode residual information for reconstruction. The same tokenizer can therefore support prefix transmission: the receiver chooses the first $k$ RVQ index layers according to channel condition and task demand. The bit cost is determined by the token grid, codebook size, and number of transmitted layers. In the present AID configuration, $k=1,2,3,4$ correspond to 2{,}560, 5{,}120, 7{,}680, and 10{,}240 bits/image, respectively.

The evidence chain of this paper has three core parts. First, RS-Token is compared with the internal plain RVQ index baseline \texttt{rvq\_baseline}, which uses the same architecture but no RemoteCLIP distillation, to test whether distillation makes L0 more semantic. Second, a layered probe checks whether later RVQ layers continue to add task semantics or mainly provide residual details. Third, reconstruction-path results evaluate whether increasing $k$ improves PSNR, LPIPS, and reconstructed-image classifier accuracy. Supplementary experiments examine whether RS-Token indices can be protected by a custom rate-1/2 systematic sparse LDPC code and how unprotected WebP/JPEG2000 bitstreams fail over BPSK channels.

The contributions are:
\begin{itemize}
    \item We propose a RemoteCLIP-guided L0 semantic allocation mechanism for a four-layer RVQ tokenizer, where only the first layer is distilled to carry remote-sensing scene semantics and later layers are left for progressive reconstruction.
    \item We formulate prefix-style RVQ index transmission for remote-sensing images. The concrete bit costs are configuration-dependent; in our AID setup, the four prefixes use 2{,}560, 5{,}120, 7{,}680, and 10{,}240 bits/image.
    \item We establish a separated evaluation protocol: $h_0$/L0 bag-of-words accuracy for the $k=1$ task path, and PSNR, LPIPS, plus a clean AID ResNet34 evaluator for the reconstruction path.
    \item We validate, on AID under AWGN and Rayleigh channels, that RemoteCLIP distillation substantially improves L0 task accuracy over plain RVQ index transmission, while additional RVQ layers improve reconstruction when the channel can support them.
\end{itemize}
\section{Related Work}
\subsection{Task-Oriented Remote-Sensing Communication}
Task-oriented semantic communication optimizes transmission for downstream decisions rather than for pixel fidelity alone. Remote-sensing applications are a natural target because scene classification, object monitoring, and emergency screening often tolerate partial visual information if the class-relevant semantics are preserved. Prior work on UAV or remote-sensing semantic communication has demonstrated that task accuracy can remain high at lower rates than conventional image reconstruction systems \cite{gao2022task}. RS-Token differs by keeping a discrete, layered representation whose first layer is deliberately aligned to a remote-sensing foundation model and whose later layers are reserved for reconstruction detail.

\subsection{Discrete Visual Tokens and Residual Quantization}
Vector-quantized autoencoders introduced discrete latent image tokens as a compact neural representation \cite{vqvae}. Residual vector quantization extends this idea by representing a latent with multiple codebooks, where later codebooks encode residual errors left by earlier ones \cite{lee2022residual}. RVQ is widely used in neural audio codecs such as SoundStream and EnCodec \cite{soundstream,encodec}, where progressive codebooks naturally define operating points. RS-Token adapts this layered-token principle to remote-sensing communication and adds a teacher-guided semantic constraint on the first layer.

\subsection{Foundation-Model Distillation}
Vision-language foundation models provide semantically structured embeddings that can supervise compact representations. RemoteCLIP is trained for remote-sensing image-text alignment and provides a domain-specific embedding space for land-cover semantics \cite{liu2024remoteclip}. Instead of distilling a generic visual representation into all latent layers, RS-Token distills RemoteCLIP only into L0-derived features. This localized distillation is important because the communication objective is hierarchical: the first layer should support robust task inference, while the remaining layers should remain available for reconstruction refinement.

\subsection{Conventional and Coded Baselines}
Conventional image codecs such as WebP and JPEG2000 are strong engineering baselines for clean image storage. In a noisy digital channel, however, an unprotected compressed bitstream may fail catastrophically because the decoder expects a syntactically valid file. Channel coding can mitigate this failure mode, but the exact code matters. In this paper, WebP/JPEG2000 results are described only as unprotected compressed bitstreams over a BPSK channel. The LDPC experiment is described separately as a rate-1/2 systematic sparse LDPC code with min-sum belief-propagation decoding; it is not a 5G NR LDPC implementation and is not used to claim comparison against a standardized cellular code.

\section{Method}
\begin{figure*}[!t]
\centering
\includegraphics[width=0.96\textwidth]{../../rstoken/figs/fig_method_aech_image2pro.png}
\caption{Overall RS-Token framework. The transmitter encodes an input remote-sensing image into four RVQ index layers. Deployment transmits a prefix of the indices, while RemoteCLIP teacher supervision and the distillation head are used only during training. The receiver can either use an L0-derived task representation or decode the received prefix for reconstruction.}
\label{fig:method_overall}
\end{figure*}

\subsection{Problem Setting}
Let $x\in\mathbb{R}^{H\times W\times3}$ be a remote-sensing image and $y$ its scene class. The transmitter maps $x$ to a hierarchy of discrete indices $z_{1:K}=\{z_1,\ldots,z_K\}$ with $K=4$. Each layer has the same spatial token grid and contributes 2{,}560 bits/image in the present implementation. A receiver may request the prefix $z_{1:k}$ according to channel quality and application demand:
\begin{equation}
    B(k)=2560k, \quad k\in\{1,2,3,4\},
\end{equation}
which gives 2{,}560, 5{,}120, 7{,}680, and 10{,}240 bits/image. The $k=1$ setting is the task-oriented low-rate mode; larger $k$ values are reconstruction-refinement modes.

\subsection{Hierarchical RVQ Encoder--Decoder}
The encoder $E_\theta$ maps an image to a continuous latent tensor $u=E_\theta(x)$. A residual vector quantizer then approximates $u$ by a sum of codebook embeddings:
\begin{align}
    r_0 &= u, \\
    z_\ell &= \arg\min_j \|r_{\ell-1}-c_{\ell,j}\|_2^2, \\
    q_\ell &= c_{\ell,z_\ell}, \\
    r_\ell &= r_{\ell-1}-q_\ell,
\end{align}
for $\ell=1,\ldots,4$. Given a prefix of $k$ layers, the decoder reconstructs
\begin{equation}
    \hat{x}_k = D_\phi\left(\sum_{\ell=1}^{k} q_\ell\right).
\end{equation}
This prefix property is what allows one trained model to support multiple transmission rates.

\subsection{RemoteCLIP Distillation on L0}
The central design choice is to align only the first quantization layer with a remote-sensing semantic teacher. Let $f_T(x)$ denote the normalized RemoteCLIP image embedding, and let $g_\psi(q_1)$ denote a projection from the L0 quantized representation to the teacher feature dimension. The distillation term is
\begin{equation}
    \mathcal{L}_{\rm distill}=1-\frac{g_\psi(q_1)^\top f_T(x)}{\|g_\psi(q_1)\|_2\|f_T(x)\|_2}.
\end{equation}
The total training objective combines reconstruction, perceptual/codebook regularization, and L0 distillation:
\begin{equation}
    \mathcal{L}=\mathcal{L}_{\rm rec}+\mathcal{L}_{\rm vq}+\lambda\mathcal{L}_{\rm distill}.
\end{equation}
The main model uses $\lambda=0.5$, selected from the earlier trade-off experiment because it balances no-channel semantic accuracy with channel robustness. The no-distillation internal baseline sets $\lambda=0$ while preserving the same architecture and training protocol.

\subsection{Channel Model and Metrics}
Index bits are modulated by BPSK and transmitted over AWGN or Rayleigh fading channels. For task-path evaluation, the receiver decodes only the L0 indices and forms an $h_0$ bag-of-words feature over L0 code usage. A linear classifier is trained/evaluated on this feature. This metric is intentionally limited to $k=1$ because it measures the semantic content of L0, not the visual quality of reconstructions produced by higher layers.

For reconstruction-path evaluation, the receiver decodes images from $z_{1:k}$ for $k=1\ldots 4$. The metrics are PSNR, LPIPS \cite{zhang2018lpips}, and top-1 accuracy of a clean AID ResNet34 classifier applied to reconstructed images. The ResNet34 evaluator is trained separately on clean AID images and reaches 96.10\% test top-1 accuracy and 95.94\% macro-F1, which makes it a stable proxy for whether reconstructions preserve scene-recognition cues. These reconstruction-path metrics, not $h_0$, support all $k=2\ldots 4$ claims.

\section{Experiments}
This section is organized around three core questions and two supplementary questions rather than the chronological order of experiment logs. First, we ask whether RemoteCLIP distillation makes L0 more semantic. Second, we test whether the later RVQ layers continue to add task semantics or mainly encode residual details. Third, we evaluate whether transmitting more RVQ layers improves reconstruction. The two supplementary experiments examine engineering boundaries: whether RS-Token indices can be protected by a conventional channel code, and how unprotected conventional compressed bitstreams fail over the same BPSK noisy channel. Each core subsection follows a problem--setting--result--conclusion structure. Conventional compression is treated as an external stress test, not as the main evidence for the proposed semantic allocation mechanism.

To avoid metric mixing, we separate the evaluation into a task path, a layered probe, and a reconstruction path. The task path uses only $h_0$/L0 bag-of-words accuracy and supports only the $k=1$ low-rate semantic claim. The layered probe uses cumulative codeword embeddings to analyze the semantic division between L0 and L1--L3. The reconstruction path uses PSNR, LPIPS, and a clean AID ResNet34 reconstructed-image classifier, and supports the $k=1\ldots4$ reconstruction claims. In other words, $h_0$/L0 bag-of-words accuracy and the layered probe are not used as evidence for multi-layer reconstruction quality, and PSNR/LPIPS alone are not used as evidence that L0 is semantic.

Experiments use the AID aerial scene dataset with 30 classes \cite{xia2017aid}. The internal tokenizer comparison uses two model families: \texttt{rvq\_baseline}, a four-layer RVQ tokenizer without RemoteCLIP distillation, and \texttt{rvq\_distill}, the same RVQ-index transmission framework with L0 RemoteCLIP distillation. The two models share the same number of RVQ layers, codebook size, and transmission protocol, so their comparison isolates the effect of RemoteCLIP distillation.

\subsection{Experimental Setup}
All experiments use a fixed train/validation/test split on AID. Training uses RandomResizedCrop(256), while evaluation uses Resize and CenterCrop(256), with image normalization matching the training scripts. RS-Token uses a four-layer RVQ tokenizer with $256\times256$ RGB input, a $16\times16$ latent token grid, latent dimension 256, and codebook size 1024. Each layer therefore contains 256 discrete symbols, each represented by $\log_2(1024)=10$ bits. The per-layer payload is 2{,}560 bits/image, giving 2{,}560, 5{,}120, 7{,}680, and 10{,}240 bits/image for $k=1,2,3,4$.

The internal comparison includes \texttt{rvq\_baseline} and \texttt{rvq\_distill}. The baseline is trained only with reconstruction and RVQ quantization losses. The distilled model has the same encoder, decoder, and four-layer quantizer, but adds the L0 RemoteCLIP distillation term defined in Section III. Both use batch size 16, AdamW, learning rate $10^{-4}$, 50 epochs, gradient clipping, and bf16 automatic mixed precision. The loss weights for L1, LPIPS, RVQ, and distillation are 1, 0.1, 1, and 0.5, respectively; LPIPS uses an AlexNet backbone. The teacher is frozen RemoteCLIP-ViT-B/32, and the DistillHead hidden dimension is 512.

For the task path, we build a 1024-dimensional bag-of-words feature $h_0$ from the received L0 indices and train a linear classifier for AID scene classification. For the reconstruction path, we decode reconstructed images from the received prefix of $k$ RVQ layers and report PSNR, LPIPS, and the accuracy of a clean AID ResNet34 classifier on reconstructed images. Unless otherwise stated, the main task-path results report mean $\pm$ standard deviation over model seeds 41/42/43, while reconstruction-path sweeps use the main configuration reported in the experiment logs. Channel experiments convert RVQ index bits to BPSK symbols, transmit them through AWGN or Rayleigh channels, and recover indices by hard decisions at the receiver.

\subsection{Does RemoteCLIP Distillation Make L0 Semantic?}
\textbf{Problem.} The central RS-Token hypothesis is that L0 should remain useful for scene-level decisions when only the lowest-rate layer is transmitted. A plain RVQ tokenizer also produces L0 indices, but because it is driven by reconstruction losses, its first codebook is not guaranteed to be discriminative for AID scene categories. This experiment asks whether RemoteCLIP distillation significantly improves L0 task fidelity under the same RVQ-index transmission framework.

\textbf{Setting.} We compare \texttt{rvq\_baseline} and \texttt{rvq\_distill}. The only intended difference is the L0 RemoteCLIP distillation term; tokenizer architecture and index transmission are otherwise kept the same. The task path uses only received L0 indices to build $h_0$/L0 bag-of-words features, followed by a linear probe for 30-class AID scene classification. This metric corresponds only to the $k=1$ low-rate task path and is not used to support reconstruction claims for $k=2\ldots4$. Table~\ref{tab:task_seed} reports mean $\pm$ standard deviation over three model seeds and includes best PSNR/LPIPS only as training-quality references, not as the main evidence for L0 semantics.

\begin{table}[!t]
\centering
\caption{Task-path stability across three model seeds. Accuracy is $h_0$/L0 bag-of-words top-1 accuracy in percent and is used only for $k=1$ task-path claims.}
\label{tab:task_seed}
\begin{tabular}{@{}lcc@{}}
\toprule
Metric & No distillation & RemoteCLIP distillation \\
\midrule
Best PSNR (dB) & $26.10\pm0.02$ & $25.89\pm0.07$ \\
Best LPIPS & $0.172\pm0.001$ & $0.175\pm0.002$ \\
No channel $h_0$ acc. & $58.23\pm1.57$ & $\mathbf{83.33\pm0.81}$ \\
AWGN $+5$ dB $h_0$ acc. & $55.73\pm0.72$ & $\mathbf{82.57\pm0.31}$ \\
AWGN $+10$ dB $h_0$ acc. & $58.20\pm1.59$ & $\mathbf{83.37\pm0.76}$ \\
Rayleigh $+5$ dB $h_0$ acc. & $28.67\pm0.65$ & $\mathbf{58.57\pm0.47}$ \\
Rayleigh $+10$ dB $h_0$ acc. & $48.63\pm1.31$ & $\mathbf{78.80\pm0.72}$ \\
\bottomrule
\end{tabular}
\end{table}

\begin{figure}[!t]
\centering
\includegraphics[width=\columnwidth]{../../rstoken/figs/fig_exp_l0_task_robustness_v3.png}
\caption{L0 task-path robustness under no-channel, AWGN, and Rayleigh conditions. The plot visualizes the same $h_0$/L0 bag-of-words task-path evidence as Table~\ref{tab:task_seed}; it is used only for $k=1$ task-path claims.}
\label{fig:l0_task_robustness}
\end{figure}

\textbf{Result.} RemoteCLIP distillation provides a large and stable L0 task gain. Without a channel, $h_0$/L0 bag-of-words accuracy improves from $58.23\pm1.57\%$ to $83.33\pm0.81\%$, a gain of 25.10 percentage points. Under AWGN $+5$ dB and $+10$ dB, the distilled model reaches $82.57\pm0.31\%$ and $83.37\pm0.76\%$, close to its no-channel performance. Rayleigh fading is more difficult, but the distilled model still outperforms the no-distillation baseline by 29.90 and 30.17 percentage points at $+5$ dB and $+10$ dB, respectively.

\textbf{Conclusion.} Table~\ref{tab:task_seed} supports the claim that RemoteCLIP distillation significantly improves the scene-level discriminability of L0 indices, making the $k=1$ low-rate task path more usable. This conclusion concerns L0 task fidelity only; it does not imply that all RVQ layers become semantic, nor does it prove reconstruction quality for $k=2\ldots4$.

\subsection{Do Later RVQ Layers Continue to Add Task Semantics?}
\textbf{Problem.} The previous subsection shows that RemoteCLIP distillation improves L0 discriminability, but this alone does not prove a clean layer specialization. A possible alternative explanation is that L1--L3 also carry substantial task semantics and L0 is only one part of a distributed semantic representation. We therefore test whether task accuracy continues to increase substantially when RVQ codewords are accumulated layer by layer.

\textbf{Setting.} For both \texttt{rvq\_baseline} and \texttt{rvq\_distill}, we construct cumulative codeword embeddings for L0, L0+L1, L0+L1+L2, and L0+L1+L2+L3, and train a linear probe on each representation for AID scene classification. This is an auxiliary mechanism analysis using the main model configuration. It analyzes semantic layer specialization; it is not a reconstruction-quality evaluation and is not used as the three-seed task-path main result.

\begin{table}[!t]
\centering
\caption{Layered linear-probe analysis on cumulative RVQ codeword embeddings. The probe is used only to analyze semantic layer specialization, not reconstruction quality.}
\label{tab:layer_probe}
\begin{tabular}{@{}lrrrr@{}}
\toprule
Layers & No dist. acc. & $\Delta$ & Distilled acc. & $\Delta$ \\
\midrule
L0 & 47.70 & -- & 86.00 & -- \\
L0+L1 & 47.20 & -0.50 & 87.70 & +1.70 \\
L0+L1+L2 & 47.70 & +0.50 & 87.90 & +0.20 \\
L0+L1+L2+L3 & 48.30 & +0.60 & 88.00 & +0.10 \\
\bottomrule
\end{tabular}
\end{table}

\textbf{Result.} For the distilled model, adding L1 to L0 increases task accuracy by only 1.70 percentage points, and adding L2 and L3 provides only another 0.30 percentage points in total. This indicates that most task semantics are already concentrated in L0. The no-distillation model stays around 47--48\% for all cumulative representations, suggesting that ordinary RVQ residual layers do not naturally form a scene-semantic hierarchy.

\textbf{Conclusion.} The layered probe supports the intended specialization: RemoteCLIP distillation mainly makes L0 scene-semantic, while later RVQ layers add little additional task semantics and are better interpreted as residual-detail layers for the reconstruction path. Therefore, claims about $k=1\ldots4$ reconstruction in the next subsection are based on reconstruction metrics, not on L0 bag-of-words accuracy.

\subsection{How Do Additional RVQ Layers Affect Reconstruction?}
\textbf{Problem.} Showing that L0 carries task semantics does not mean that L0 alone is sufficient for image reconstruction. This experiment asks whether transmitting additional RVQ layers improves image quality and reconstructed-image classification accuracy on the reconstruction path.

\textbf{Setting.} Table~\ref{tab:recon_path} reports reconstruction-path results for the RemoteCLIP-distilled tokenizer. The $h_0$/L0 metric is not used in this subsection. Claims about $k=2\ldots4$ are supported only by PSNR, LPIPS, and clean AID ResNet34 reconstructed-image classifier accuracy. The full $k=1\ldots4$ sweep is available for the no-channel and AWGN $+5$ dB settings; AWGN $+10$ dB and Rayleigh settings report key $k=4$ operating points to characterize high-SNR or fading behavior.

\begin{table*}[!t]
\centering
\caption{Reconstruction path for the RemoteCLIP-distilled tokenizer. The $h_0$/L0 metric is not used here; $k=2\ldots4$ claims are supported by reconstructed-image metrics.}
\label{tab:recon_path}
\begin{tabular}{@{}llrrrrr@{}}
\toprule
Channel & SNR & $k$ & Bits/image & PSNR (dB) & LPIPS $\downarrow$ & Recon. cls. acc. (\%) \\
\midrule
None & -- & 1 & 2{,}560 & 22.95 & 0.284 & 70.3 \\
None & -- & 2 & 5{,}120 & 24.77 & 0.205 & 84.3 \\
None & -- & 3 & 7{,}680 & 25.57 & 0.183 & 86.4 \\
None & -- & 4 & 10{,}240 & 25.92 & 0.175 & 86.9 \\
\midrule
AWGN & $+5$ dB & 1 & 2{,}560 & 21.97 & 0.312 & 67.3 \\
AWGN & $+5$ dB & 2 & 5{,}120 & 23.24 & 0.244 & 81.6 \\
AWGN & $+5$ dB & 3 & 7{,}680 & 23.76 & 0.223 & 84.1 \\
AWGN & $+5$ dB & 4 & 10{,}240 & 23.96 & 0.217 & 84.3 \\
AWGN & $+10$ dB & 4 & 10{,}240 & 25.92 & 0.175 & 86.9 \\
Rayleigh & $+5$ dB & 4 & 10{,}240 & 16.96 & 0.467 & 21.0 \\
Rayleigh & $+10$ dB & 4 & 10{,}240 & 20.60 & 0.323 & 66.6 \\
\bottomrule
\end{tabular}
\end{table*}

\begin{figure}[!t]
\centering
\includegraphics[width=\columnwidth]{../../rstoken/figs/fig_exp_progressive_reconstruction_v3.png}
\caption{Progressive reconstruction as additional RVQ layers are transmitted. The figure corresponds to the reconstruction-path metrics in Table~\ref{tab:recon_path}; it does not use $h_0$/L0 bag-of-words task metrics.}
\label{fig:progressive_reconstruction}
\end{figure}

\textbf{Result.} Without a channel, increasing $k$ from 1 to 4 improves PSNR from 22.95 dB to 25.92 dB, reduces LPIPS from 0.284 to 0.175, and raises reconstructed-image classifier accuracy from 70.3\% to 86.9\%. AWGN $+5$ dB shows the same direction: reconstructed-image classifier accuracy increases from 67.3\% at $k=1$ to 84.3\% at $k=4$. The AWGN $+10$ dB $k=4$ result is nearly identical to the no-channel $k=4$ result. Rayleigh $+5$ dB is a more difficult fading stress case: even at $k=4$, reconstructed-image classifier accuracy is only 21.0\% in the unprotected reconstruction-path sweep.

\textbf{Conclusion.} Table~\ref{tab:recon_path} supports the conclusion that additional RVQ layers progressively improve the reconstruction path in the full no-channel and AWGN $+5$ dB sweeps. Under strong Rayleigh fading, reconstruction still degrades substantially. This conclusion is supported only by reconstruction-path metrics and does not rely on $h_0$/L0 bag-of-words accuracy.

\subsection{Can RS-Token Indices Be Protected by Channel Coding?}
\textbf{Problem.} The preceding experiments send RS-Token index bits directly over BPSK noisy channels to observe task and reconstruction degradation under index errors. Practical digital systems often use error-correcting codes. This supplementary experiment therefore asks whether RS-Token indices can be treated as an ordinary digital bitstream and protected by channel coding.

\textbf{Setting.} We protect the RS-Token bitstream with a rate-1/2 systematic sparse LDPC code and use min-sum belief-propagation decoding at the receiver. This is a custom sparse LDPC implementation, not a 5G NR LDPC code. Because the coding rate is 1/2, transmitted bits double relative to the uncoded source indices: $k=1$ changes from 2{,}560 source bits/image to 5{,}120 transmitted bits/image, and $k=4$ changes from 10{,}240 source bits/image to 20{,}480 transmitted bits/image.

\begin{table}[!t]
\centering
\caption{Selected rate-1/2 systematic sparse LDPC results for RS-Token, averaged over five channel seeds. This table is a channel-coding extension and should not be described as 5G NR LDPC.}
\label{tab:ldpc_rstoken}
\begin{tabular}{@{}llrrrr@{}}
\toprule
Channel & SNR & $k$ & Tx bits & Recon. cls. acc. (\%) & $h_0$ acc. (\%) \\
\midrule
AWGN & $+5$ dB & 1 & 5{,}120 & $70.1\pm0.2$ & $82.4\pm0.2$ \\
AWGN & $+10$ dB & 1 & 5{,}120 & $70.3\pm0.0$ & $82.4\pm0.0$ \\
Rayleigh & $+5$ dB & 1 & 5{,}120 & $47.2\pm0.6$ & $75.0\pm0.8$ \\
Rayleigh & $+10$ dB & 1 & 5{,}120 & $69.1\pm0.6$ & $82.6\pm0.4$ \\
AWGN & $+10$ dB & 4 & 20{,}480 & $86.9\pm0.0$ & -- \\
Rayleigh & $+10$ dB & 4 & 20{,}480 & $86.1\pm0.1$ & -- \\
\bottomrule
\end{tabular}
\end{table}

\textbf{Result.} Under AWGN $+5$ dB and $+10$ dB, the $k=1$ $h_0$ accuracy remains around 82.4\%, and reconstructed-image classifier accuracy is about 70\%. Under Rayleigh $+10$ dB, $k=1$ $h_0$ accuracy reaches 82.6\%, and $k=4$ reconstructed-image classifier accuracy reaches 86.1\%. Rayleigh $+5$ dB remains difficult: the $k=1$ reconstructed-image classifier accuracy drops to $47.2\pm0.6\%$, indicating that channel coding improves robustness but does not fully remove the effect of strong fading.

\textbf{Conclusion.} This experiment supports a compatibility claim: RS-Token indices can be protected as ordinary digital bitstreams by an error-correcting code. It is not a new channel-coding method and is not a 5G NR LDPC comparison. The main evidence for the paper remains the L0 task-path experiment, the layered specialization probe, and the multi-layer reconstruction-path experiment.

\subsection{How Do Unprotected Conventional Bitstreams Fail over BPSK Channels?}
\textbf{Problem.} The main internal comparison for RS-Token is \texttt{rvq\_baseline}, because it shares the same RVQ-index framework and directly tests the effect of RemoteCLIP distillation. WebP and JPEG2000 are external conventional compression baselines. They are not architectural ablations and should not be treated as a perfectly fair semantic-tokenizer comparison. This supplementary experiment asks a narrower question: how do conventional compressed files fail when their bitstreams are sent unprotected over BPSK noisy channels?

\textbf{Setting.} WebP and JPEG2000 are encoded at target bit budgets of 2{,}560, 5{,}120, and 10{,}240 bits/image, aligned with RS-Token $k=1$, $k=2$, and $k=4$. The conventional-compression experiment has no target-bits setting at 7{,}680 bits/image and therefore does not cover the RS-Token $k=3$ operating point. The compressed bitstreams are sent directly over BPSK channels without LDPC or any other error-correcting code. Actual bits denote the mean generated file size. Decode failure is the fraction of images that cannot be decoded by the receiver. All-image classifier accuracy counts failed decodes as classification errors, while valid PSNR is computed only on successfully decoded images.

\begin{table*}[!t]
\centering
\caption{External conventional-compression stress test. These are unprotected WebP/JPEG2000 compressed bitstreams over BPSK channels; they are not LDPC-protected. Target bits align with RS-Token $k=1,2,4$ only.}
\label{tab:classic_unprotected}
\begin{tabular}{@{}llrrrrr@{}}
\toprule
Method & Channel/SNR & Target bits & Actual bits (mean) & Decode failure & All-image cls. acc. (\%) & Valid PSNR (dB) \\
\midrule
WebP & None & 2{,}560 & 14{,}164 & 0.000 & 61.8 & 26.56 \\
WebP & AWGN $+10$ dB & 2{,}560 & 14{,}164 & 0.039 & 58.8 & 26.51 \\
WebP & Rayleigh $+10$ dB & 2{,}560 & 14{,}164 & 0.999 & 0.0 & 10.52 \\
WebP & None & 5{,}120 & 14{,}301 & 0.000 & 62.6 & 26.65 \\
WebP & AWGN $+10$ dB & 5{,}120 & 14{,}301 & 0.045 & 59.9 & 26.63 \\
WebP & Rayleigh $+10$ dB & 5{,}120 & 14{,}301 & 1.000 & 0.0 & -- \\
WebP & None & 10{,}240 & 15{,}410 & 0.000 & 67.5 & 27.04 \\
WebP & AWGN $+10$ dB & 10{,}240 & 15{,}410 & 0.041 & 64.5 & 26.97 \\
WebP & Rayleigh $+10$ dB & 10{,}240 & 15{,}410 & 1.000 & 0.0 & -- \\
\midrule
JPEG2000 & None & 2{,}560 & 2{,}663 & 0.000 & 6.5 & 20.19 \\
JPEG2000 & AWGN $+10$ dB & 2{,}560 & 2{,}663 & 0.003 & 6.6 & 20.11 \\
JPEG2000 & Rayleigh $+10$ dB & 2{,}560 & 2{,}663 & 1.000 & 0.0 & -- \\
JPEG2000 & None & 5{,}120 & 5{,}192 & 0.000 & 11.0 & 22.35 \\
JPEG2000 & AWGN $+10$ dB & 5{,}120 & 5{,}192 & 0.001 & 10.9 & 22.27 \\
JPEG2000 & Rayleigh $+10$ dB & 5{,}120 & 5{,}192 & 1.000 & 0.0 & -- \\
JPEG2000 & None & 10{,}240 & 10{,}246 & 0.000 & 39.6 & 23.97 \\
JPEG2000 & AWGN $+10$ dB & 10{,}240 & 10{,}246 & 0.004 & 38.9 & 23.90 \\
JPEG2000 & Rayleigh $+10$ dB & 10{,}240 & 10{,}246 & 1.000 & 0.0 & -- \\
\bottomrule
\end{tabular}
\end{table*}

\textbf{Result.} Conventional compressed bitstreams are useful clean-reference baselines, but unprotected bitstreams show clear file-level fragility under noisy channels. WebP has higher clean reconstructed-image classifier accuracy, but its actual bit counts are far above the requested target bits in this implementation, so it is not a strict rate-matched comparator. Under Rayleigh $+10$ dB, WebP decoding fails for essentially all images. JPEG2000 more closely matches the requested target bits, but its clean classifier accuracy is low at small bit budgets and it also suffers many decode failures under AWGN $+10$ dB and Rayleigh $+10$ dB.

\textbf{Conclusion.} This section supports only a limited conclusion: unprotected conventional compressed bitstreams are prone to file-level failure over BPSK noisy channels, whereas RS-Token indices are easier to separate into task and reconstruction representations that can be used layer by layer. This experiment is not a complete comparison against an engineered ``codec + strong channel code + retransmission'' system, and it does not replace \texttt{rvq\_baseline} as the core baseline for testing the contribution of RemoteCLIP distillation.

\section{Discussion}
The experiments support three core conclusions and two supplementary conclusions, each tied to a distinct metric family. First, the semantic-layer conclusion is a task-path conclusion: across three model seeds, L0 bag-of-words accuracy improves from $58.23\pm1.57\%$ to $83.33\pm0.81\%$. Second, the layered probe shows that later RVQ layers do not add substantial task semantics, supporting the division that L0 is task-biased while L1--L3 mainly encode residual details. Third, the progressive-reconstruction conclusion is a reconstruction-path conclusion: in the full no-channel and AWGN $+5$ dB $k=1\ldots4$ sweeps, PSNR, LPIPS, and reconstructed-image classifier accuracy generally improve as $k$ increases. The LDPC extension supports only a compatibility conclusion: RS-Token indices can be combined with a conventional error-correcting code. The WebP/JPEG2000 results support only an external stress-test conclusion: unprotected conventional compressed bitstreams can fail to decode over BPSK noisy channels and should not be described as a complete codec-plus-channel-code comparison.

RemoteCLIP helps the first layer because it provides a domain-specific semantic geometry for remote-sensing images. Distilling this geometry into L0 makes the first codebook layer more discriminative for aerial scene categories than a reconstruction-only RVQ layer. This does not mean that all layers become semantic, nor that L0 alone is sufficient for high-fidelity reconstruction. A more precise interpretation is layer specialization: L0 is biased toward task semantics, while L1--L3 improve pixel and perceptual detail.

Teacher ablation further bounds the RemoteCLIP claim, but those auxiliary numbers are not part of the v0.4 evidence tables used for the main results. The safe interpretation is that the main tables support L0 vision-language distillation under the RemoteCLIP teacher, while teacher-specific superiority should be treated as an auxiliary observation rather than a primary v0.4 claim.

The distillation-weight sweep also suggests a trade-off between semantic alignment and channel robustness, but its auxiliary numbers are not part of the v0.4 table set. The selected weight should therefore be described conservatively as the main configuration used in the reported v0.4 experiments, rather than as a fully optimized universal choice.

From a system perspective, RS-Token is best viewed as an application-layer representation above physical-layer adaptation, not as a replacement for HARQ, adaptive modulation and coding, or LDPC. When the link is poor or the receiver only needs an alert, the system can transmit L0 and use the task path. When the link improves or human inspection is required, it can request L1--L3 to improve reconstruction. The LDPC experiment shows that the index representation can be protected by an error-correcting code, but jointly optimizing layer selection, retransmission, modulation, and channel coding remains future work.

Several limitations remain. The experiments use AID scene classification rather than object detection or semantic segmentation, so the current task-path evidence is strongest for image-level land-cover semantics. Reconstruction-path sweeps are reported for the main model seed, while task-path claims are supported by three model seeds. We do not provide a strict rate-matched comparison against a complete conventional codec-plus-channel-code system. The LDPC experiment uses a custom sparse LDPC code with min-sum decoding and should not be presented as a 5G NR LDPC comparison.

\section{Conclusion}
This paper presented RS-Token, a hierarchical RVQ tokenizer for remote-sensing communication. By distilling RemoteCLIP into the first quantization layer, RS-Token makes L0 more task-discriminative, while later layers progressively refine reconstruction. The evaluation separates L0 task-path evidence from multi-layer reconstruction-path evidence and uses plain RVQ index transmission as the internal baseline for testing the contribution of distillation. Supplementary unprotected WebP/JPEG2000 and custom rate-1/2 systematic sparse LDPC experiments clarify the external stress-test and channel-coding compatibility boundaries. Overall, the results support a conservative but useful claim: RemoteCLIP-guided layer specialization can provide low-rate task fidelity and progressive reconstruction for channel-robust remote-sensing image communication.
\bibliographystyle{IEEEtran}
\bibliography{rs_token}

\end{document}








